\scanpage{128}
\unsure{Posons|poser}
\[
B=H\mathbin{\amalg}_k H'=H\oplus H'/\deleted{\unsure{k1_H+k1_{H'}|k(1_H+1_{H'})}}\inserted{k(1_H\ominus1_{H'})}
\]
% slanted
\margin{Mais la structure d'alg. de $H$ résulte déjà de sa structure de module quadratique, ``muni de $1_H$ tel que $\Delta_H(1)=1$'', et itou pour $H'$ !} Ainsi $H$ et $H'$ se réalisent comme sous-fibrés vect. de rg 2 d'un $B$, et
\[
H\cap H'=k1_B\qquad H+H'=B
\]
où $1_B=1_H=1_{H'}$ est l'image commune de $1_H$ et $1_{H'}$ dans $B$.

On considère aussi
\[
P=H\otimes_k H'
\]
et les plong\textsuperscript{ts}
\diagram{\begin{tikzcd}[ampersand replacement=\&,column sep=small,row sep=small]
\&P\&\&[-1em] \\ H\arrow[ur,"i_0"]\& \& H'\arrow[ul,"j_0"']\& {\begin{gathered}i(u)=u\otimes1\\ j(v)=1\otimes v\end{gathered}}
\end{tikzcd}}
On peut alors réaliser $B$ comme sous-fibré vectoriel de $P$, \deleted{\illegible{1}}\inserted{car} $i,j$ [\deleted{\unsure{le|lequel}} qui ``coïncident sur $k$'' i.e. $i(\lambda1_H)=j(\lambda1_{H'})=\lambda\underbrace{1_H\otimes1_{H'}}_{\text{``est'' }1_P}\bigr]$ se factorisent de façon unique par $B\longrightarrow P$
\diagram{\begin{tikzcd}[ampersand replacement=\&,column sep=small,row sep=small]
\&P\&\\ \&B\arrow[u]\arrow[dl]\arrow[dr]\&\\ H\arrow[uur,"i_0"]\&\&H'\arrow[uul,"j_0"']
\end{tikzcd}}
% slanted
\margin{Attention, \unsure{Priori pas|$B$ n'a pas|$P$ n'a pas} de structure d'algèbre !} \margin{NB on a $B\xrightarrow{T_B}k$ induisant $T_H,T_{H'}$}
% slanted
\margin{(puisque $T_H(1_H)=T_{H'}(1_{H'})=2$)}
% Faint mirror-image typescript shows through the sheet.
% archive sheet number: 127
